\section{Do the highest-prevalence patients shape the factors?}
\label{sec:shape}

\begin{quote}
``We discussed quite a bit on the observation that patients with more malignant cells (P02, P03, P04) disproportionally shape the factors.''\\[2pt]
TODO: ``Re-weight cells per patient in Stage 1 so each patient contributes equally to DR \ldots{} then re-examine PCA and FA factors to see if the imbalance problem persists.''
\end{quote}

{\small\color{inkmuted}\emph{Reading guide.} What a leverage test is: the box ``The leverage test, step by step''; what it found on allograft rejection and apoptosis, the run you started (D26): section~\ref{glance:02b}. The metrics: the boxes ``Score energy, intuitively'' and ``The leverage test''. Figure~\ref{fig:02a} keeps cells, labels and score energy in separate panels; Figure~\ref{fig:02b} has its axis in percent and a reading of one row. How re-weighting works and what its percentages mean: the box ``How patient-equal re-weighting works''; your choice of target (2026-10-10): the options list. What the section implies for healthy-cell downsampling: the block after the re-weighting options.\par}

\begin{answerbox}
\textbf{Not more than other patients do, on the three panels tested (PCA and FA).}
\par\smallskip\noindent \emph{IFN-$\alpha$:} refitting Stage~1 without P02+P03+P04 leaves the malignant-vs-normal axis almost unchanged (under 1\% of its variance lost; matched $|r| \geq 0.988$) and not beyond random drops of the same number of cells and malignant cells (above all 20 in 3 of 6 method $\times$ $K$ cells; the rule needs 4, and on this panel it is applied post hoc).
\par\smallskip\noindent \emph{Allograft rejection and apoptosis} (D26, rule fixed before the run; section~\ref{glance:02b}): leaving them out moves the axis beyond those random drops (5 of 6 and 6 of 6 cells, mostly P03), but never more than leaving out other patients does; P01 moves it most. By the rule, the hypothesis that they shape the factors more than other patients do is not supported.
\par\smallskip\noindent \emph{Score space:} they hold about their 8.2\% cell share of factor variance on all 41 panels (median ratio 0.96; 1.04 after adjusting for cell type). \emph{Where they dominate:} Stage~2, with 68.3\% of the malignant labels. \emph{Not tested:} FactoSig bases (the atlas winners on allograft and apoptosis) and the other 38 panels.
\par\smallskip\noindent \emph{Decision it informed:} what re-weighting should target (W2.3). On the three panels tested (PCA and FA), the imbalance specific to these patients sits mainly in the Stage~2 labels: in the Stage~1 geometry it goes beyond composition on two panels (mostly P03) but, on the primary axis, never beyond other patients, and the largest Stage~1 leverage is P01's. You chose the Stage~2 labels on 2026-10-10, for the runs with the new patient data.
\end{answerbox}

\subsection*{What we did}

P03, P02 and P04 are the \emph{highest-prevalence} patients (32.6\%, 14.4\% and 9.0\% malignant), not the ones with the most malignant cells (P09 has more than P02 or P04); by total cells they are among the smallest. The stored Stage~1 bases are \emph{transductive}: fit once, without labels, on all 60,118 cells of the 13 patients, including any patient later held out and the four patients with no malignant cell at all. ``Shaping the factors'' can mean two things, and we measured both:
\begin{enumerate}[label=(\arabic*)]
\item \textbf{Where the variance sits.} On the fixed stored basis, how much of the factor-score variance sits in each patient's cells, compared with its share of the cells (``score energy'', next box; all 320 stored PCA/FA bases). D28 adds a cell-type adjustment: is a patient's excess just its cell-type mix?
\item \textbf{Leverage.} An in-silico perturbation: remove the patients from the Stage~1 fit, refit, and measure how far the malignant axis moves, against random removals of the same size and malignant composition (box ``The leverage test, step by step''; IFN-$\alpha$ here, allograft rejection and apoptosis in section~\ref{glance:02b}; PCA and FA at $K$ = 10, 20, 40).
\end{enumerate}
Only (2) tests whether a patient \emph{made} a factor; (1) is a cheap screen over all panels.

\begin{explainbox}{Score energy, intuitively}
\emph{Picture.} Each cell is a point in the $K$-dimensional space of its factor scores. Its \emph{score energy} is its squared distance from the cohort centre, $e_i = \sum_f (z_{if} - \bar z_f)^2$, on the stored (raw, not standardized) scores, so a high-variance factor counts more. Summed over all cells, the energy is $n$ times the total variance of the $K$ factor scores: for PCA exactly the quantity the factors are chosen to maximise.

\emph{Share and ratio.} A patient's \emph{energy share} is $\sum_{i \in p} e_i / \sum_i e_i$: the fraction of the captured variance that sits in its cells. Its \emph{ratio} is energy share $\div$ cell share, which equals the patient's mean energy per cell divided by the cohort's. A ratio of 1 means its cells are as spread out as an average cell; 1.11 means 11\% more squared distance per cell.

\emph{Worked example} (Figure~\ref{fig:02a}b). On IFN-$\alpha$ PCA40, P02+P03+P04 are 8.2\% of the cells and hold 9.1\% of the energy: $9.1 / 8.2 = 1.11$. On FA40 they hold 8.4\%: 1.02.

\emph{What it is not.} A share is an allocation on the fixed basis, not proof that those cells pulled the factors (that is the leverage test). It also ignores labels: the malignant-cell share in Figure~\ref{fig:02a}a is a different quantity, the Stage~2 label imbalance. With every factor standardized first (each counts equally), the 320-base median ratio of P02+P03+P04 is 1.00 instead of 0.96. Energy is summed over all of a patient's cells, malignant ones included; restricted to normal cells, the group's IFN-$\alpha$ ratio is still 1.11 (PCA40; FA40 1.03) and its 320-base median 0.99, so the excess is not carried by their malignant cells. \src{2026-10-09\_d28-celltype-score-shares}
\end{explainbox}

\subsection*{What we found}

\textbf{The labels are concentrated; the cells and the score energy are not} (Figure~\ref{fig:02a}).
Stage~1 sees cells, and the cell-heavy patients are P09 (28.5\%) and P01 (19.7\%). P02+P03+P04 are 4,917 of 60,118 cells (8.2\%) but 1,071 of 1,569 malignant cells (68.3\%; P03 alone 48.3\%). The four patients with no malignant cells (P08, P10, P11, P12) hold 17.7\% of the cells in the fit. On IFN-$\alpha$ at $K = 40$ the group's ratio is 1.11 (PCA) and 1.02 (FA); per patient, the most over-represented on this panel are P02 (1.43 PCA40, 1.29 FA40) and P07 (1.31, 1.44). PCA40's 1.11 is the third highest of the 320 bases, so IFN-$\alpha$ is a favourable case for the hypothesis and still near parity. Over all 320 bases the group's ratio has median 0.96 (range 0.58--1.14; above 1 in 75/320); per patient the medians are P01 1.24 (above 1 in 279/320), P02 1.06, P03 0.91 and P04 0.89, and the group's ratio creeps up with $K$ (PCA medians 0.92, 0.94, 0.98, 1.02 at $K$ = 5, 10, 20, 40) because its per-factor excess sits on low-variance trailing factors. \src{2026-10-08\_d16-factor-patient-matrix}

\begin{figure}[H]
\centering
\includegraphics[width=\linewidth]{figures/fig02a_who_dominates.pdf}
\caption{Two different imbalances, kept apart. (a) Each patient's share of the 60,118 cells (light; also its share of the unweighted Stage~1 fit) and of the 1,569 malignant cells (dark; its share of the Stage~2 positive labels). Dashed line: the share every patient would get under patient-equal weighting (1/13); over the P02+P03+P04 column, what the three get together (3/13). (b) Score-energy ratio on IFN-$\alpha$ at $K = 40$ (box ``Score energy, intuitively''), PCA40 blue, FA40 orange; a stem above 1 means the patient's cells hold more factor variance than their head count. The P02+P03+P04 values, 1.11 and 1.02, are labelled. (c) The same combined ratio on every stored PCA/FA basis; the ringed dots are IFN-$\alpha$ $K = 40$. Look at P03 in (a): 3.9\% of the cells, 48.3\% of the labels; then at its ratio in (b): 1.00 and 0.91.}
\label{fig:02a}
\end{figure}

\begin{explainbox}{The leverage test, step by step}
An in-silico perturbation, run on IFN-$\alpha$ (D17, this section) and on allograft rejection and apoptosis (D26, section~\ref{glance:02b}), for PCA and FA at $K$ = 10, 20 and 40:
\begin{enumerate}[label=\arabic*.,leftmargin=1.4em]
\item \emph{Reference axis.} On the stored basis, the factor with the largest standardized malignant--normal mean gap: PC8 (PCA) and FA7 (FA) at every $K$ on IFN-$\alpha$ (PC8 and FA8 for allograft rejection, PC6 and FA6 for apoptosis). Under balanced L1 it is the factor a sparse probe takes first (largest entry margin).
\item \emph{Perturb.} Refit Stage~1 without the patient set (genes re-standardized on the retained cells; same solver and seed).
\item \emph{Measure.} Over the retained cells, regress the original reference-axis scores on all $K$ refit factors (least squares with intercept). $R^2$ is the share of the axis's variance the refit basis still reproduces, and \textbf{axis loss} $= 1 - R^2$. Scale: 0 means fully recoverable; $10^{-3}$ means 0.1\% of the axis's variance is lost (best linear reconstruction correlates $R = \sqrt{1 - 10^{-3}} = 0.9995$ with the original); $10^{-2}$ means 1\% ($R = 0.995$). \textbf{Matched $|r|$} is the correlation with the single best-matching refit factor (one-to-one matching). Small axis loss with a low matched $|r|$ means the axis survives but is spread over several refit factors (rotated).
\item \emph{Compare.} Any refit on fewer cells moves the basis a little, so steps 2--3 are repeated for random drops. \emph{Size-matched} controls (3 draws) drop as many random cells. \emph{Composition-matched} controls (20 draws) drop as many cells \emph{and} as many malignant cells, at random from the whole cohort. Why both: the reference axis is the malignant-vs-normal direction, so removing malignant cells moves it whoever they belong to (dropping all 1,569 malignant cells moves the PCA40 axis 17.1$\times$ as much as dropping 1,569 random cells). A random draw of 4,917 cells holds about 128 malignant cells; P02+P03+P04 hold 1,071. Only the composition-matched draws separate ``these patients' cells'' from ``this many malignant cells''.
\item \emph{Summarise.} Ratio = observed axis loss $\div$ the mean of the 20 composition-matched controls; rank $p = (1 + \#\{\text{controls} \geq \text{observed}\})/21$, whose floor 0.048 means ``above all 20''. Descriptive and uncorrected.
\end{enumerate}
\emph{Reading one row} (Figure~\ref{fig:02b}a, P02+P03+P04 at PCA40): axis loss 0.0015, i.e.\ 0.15\% of PC8's variance (matched $|r|$ 0.995). Its 20 composition-matched controls lose 0.00054--0.00101 (mean 0.00079), so it is above all 20: rank $p$ 0.048, ratio 1.93. Against the size-matched mean (0.00030) it would have looked like 5.2$\times$.

\emph{Decision rule}, fixed for the allograft/apoptosis run (D26) before it ran: leverage beyond composition means above all 20 controls in at least 4 of the 6 method $\times$ $K$ cells; support for ``the highest-prevalence patients shape the factors'' also needs their ratio to exceed that of every contrast set (P01, P06, P09, P01+P09) in at least 4 cells. Both conditions were written after the IFN-$\alpha$ results below were known, so on IFN-$\alpha$ the rule is a post-hoc reading, not a pre-registered test; for allograft and apoptosis it is pre-specified.
\end{explainbox}

\textbf{On IFN-$\alpha$, leaving them out barely moves the axis; most of what moves is malignant composition} (Figure~\ref{fig:02b}).
Without P02+P03+P04 the IFN-$\alpha$ axis is essentially unchanged: PCA40 axis loss 0.0015 (matched $|r|$ 0.9953), FA40 0.00062 (0.9958); across method and $K$ the matched $|r|$ is 0.9885--0.9969 and the largest loss is 0.0069 (PCA10, 0.7\% of the axis). Size-matched controls had made this look like a 3.1--5.2$\times$ excess. Against composition-matched controls the FA excess disappears (1.00, 0.88, 0.36$\times$ at $K$ = 10, 20, 40; rank $p$ 0.48, 0.86, 0.19) and the PCA excess shrinks to 1.42--1.94$\times$ (rank $p$ 0.048 each). That is 3 of 6 cells above all controls, below the rule's 4, and in 0 of 6 cells is the group's ratio larger than every contrast set's. Panel~(c) shows why the second condition is needed: clearing composition-matched controls is common for any real patient (P01 and P09 in 6/6 cells, P02 and P06 in 5/6, P03 in 4/6, P04 in 2/6). The group clears fewer cells than P02 alone because its controls drop 1,071 random malignant cells, about two thirds of them from P02--P04 themselves, a stricter yardstick than a single patient's. A plausible reason: one patient's cells are a coherent block (one sample, one processing batch, one cell-type mix), while the controls are scattered across the cohort. P01 is the outlier (PCA 15.0--23.8$\times$, FA 3.74--21.0$\times$ its controls): without it the PCA40 axis loss is 0.0118 and PC8's best single match falls to $|r|$ 0.58, still in the subspace but rotated. \src{2026-10-08\_d17-basis-variant-pilot; 2026-10-08\_d17-composition-matched-controls}

\begin{figure}[H]
\centering
\includegraphics[width=\linewidth]{figures/fig02b_axis_leverage.pdf}
\caption{Leave-out refits of the IFN-$\alpha$ basis. (a, b) $K = 40$. Each row is one set removed from the Stage~1 fit (its cells and malignant cells on the left). Coloured dot: axis loss of PC8 (a) or FA7 (b) after refitting without that set. Bottom axis: axis loss as a percentage of the reference axis's variance that the refit factors cannot reproduce (log); top axis: the same value as $R = \sqrt{1 - \text{loss}}$. Gray ticks above the dot: the 3 size-matched random drops; dark dots below it: the 20 composition-matched random drops. $|r|$: matched correlation for the rows discussed in the text. (c) All $K$: axis loss divided by the mean of its 20 composition-matched controls (triangle, square, circle = $K$ 10, 20, 40); filled = above all 20 (rank $p$ 0.048), open = inside the control range; right column: cells above all 20, out of 6. Look at the P02+P03+P04 row: in (a) the blue dot sits just right of its control cloud (0.15\% vs 0.05--0.10\%); in (c) PCA is filled and FA is open, 3/6. For FA40, 2--10 of each set's 20 composition-matched drops fall in a second mode at 1.1--1.8\% (drawn as a separate segment), so FA40 dots there are within control noise.}
\label{fig:02b}
\end{figure}

\textbf{What lies behind P01? Its cell-type mix matches its score-energy excess, not its leverage} (Figure~\ref{fig:02c}).
D28 gives every cell the average energy of its predicted cell type (13 types) and adds these up per patient: the energy its cell-type mix alone predicts. The \emph{adjusted} ratio is observed $\div$ expected, an indirect standardization like a standardized mortality ratio; 1 means the mix fully accounts for the patient's ratio. Over the 320 bases, P01 goes from 1.24 (observed) to 1.21 (expected from its mix) and 1.03 (adjusted); the mix accounts for a median 90\% of its log departure from 1. The mechanism is in panel~(b): ncMono cells carry 3.25$\times$ the cohort's energy per cell, pDC2 1.26$\times$ and cDC 1.18$\times$, while HSC/MPP2 carry 0.65$\times$; P01's cells are 20.5\% pDC2, 14.2\% cDC and 8.4\% ncMono. P02+P03+P04 have a low-energy mix (16.3\% HSC/MPP2 at 0.65$\times$, 11.4\% MEP-EP and 10.6\% NP1 at 0.88$\times$ and 0.87$\times$, almost no ncMono or pDC2; their largest type, CLP at 29.1\%, sits at the cohort average): 0.96 observed, 0.93 expected, 1.04 adjusted (above 1 in 230/320). The IFN-$\alpha$ 1.11 is not the mix: PCA40 1.11 observed, 1.00 expected, 1.11 adjusted; FA40 1.02, 0.95, 1.08. Nor is P02's excess (1.06, 0.96, 1.11 over the 320 bases), or P07's on IFN-$\alpha$ (PCA40 1.31, 0.85, 1.54). The adjustment depends on the reference for P01, which holds 93\% of the cohort's ncMono cells: with the other patients' cells as the reference, P01's adjusted ratio is 1.08 and the mix accounts for 67\%.

``Accounts for'' is arithmetic here: the mix predicts most of P01's \emph{score energy}, which does not show that cell type causes it (the labels are transferred calls, and batch, sampling time and patient biology stay in the comparison). And it says nothing about P01's \emph{leverage}: on IFN-$\alpha$ PCA40, P01's ratio is 0.98 (adjusted 0.92), yet removing it moves PC8 15$\times$ as much as its composition-matched controls, the largest ratio of any set. One candidate, untested: leaving P01 out removes 93\% of the cohort's ncMono cells, the highest-energy cell type, almost a whole cell population. A leave-out control matched on cell type would test it. \src{2026-10-09\_d28-celltype-score-shares; 2026-10-08\_d17-composition-matched-controls}

\begin{figure}[H]
\centering
\includegraphics[width=\linewidth]{figures/fig02c_celltype_adjusted.pdf}
\caption{Is a patient's score-energy ratio its cell-type mix? (a) Open: the observed ratio (the quantity of Figure~\ref{fig:02a}b), here as the median (bar: interquartile range) over the 320 stored PCA/FA bases. Filled: the same ratio after the cell-type adjustment (observed $\div$ what the patient's predicted cell-type mix alone predicts). Right: the P02+P03+P04 values on IFN-$\alpha$ PCA40 (blue) and FA40 (orange). (b) Energy per cell of each predicted cell type relative to the cohort's (log; median and IQR over the 320 bases), with each type's percentage of P01's and of P02+P03+P04's cells underneath (bold: at least 10\%). Look at P01 in (a): its ratio falls from 1.24 to 1.03 once its pDC2, cDC and ncMono cells are accounted for; P02's rises from 1.06 to 1.11. Raw-score energy, cohort cell-type averages, all cells.}
\label{fig:02c}
\end{figure}

\textbf{Allograft rejection and apoptosis: beyond composition, mostly P03, but not beyond the other patients} (section~\ref{glance:02b}). The same test on the two panels where DR gains most over direct genes, with the rule fixed before the run (D26): without P02+P03+P04 the primary axis moves more than all 20 composition-matched drops in 5 of 6 cells (allograft) and 6 of 6 (apoptosis), at 1.31--4.15$\times$ and 2.01--3.98$\times$ their mean, mostly through P03; but the group's ratio beats every contrast set's in 0 of 6 primary-axis cells on each panel (1 of 6 on the apoptosis probe axis), and P01 moves the same axes 21--141$\times$ and 48--128$\times$ its controls. Section~\ref{glance:02b} has the design, the figures, the absolute sizes and what the result means for re-weighting and downsampling. \src{2026-10-10\_d26-leverage-allograft-apoptosis-done}

% Keep this heading with its explainer box: a new page only when fewer than twelve lines are left (the box
% does not start in a short remainder, which left the heading alone at a page foot). Same macro as report.tex;
% \providecommand keeps the standalone compile working.
\makeatletter
\providecommand{\softneedspace}[1]{\par\begingroup\dimen@ii\pagegoal\advance\dimen@ii-\pagetotal
  \ifdim #1>\dimen@ii\ifdim\dimen@ii>\z@\newpage\fi\fi\endgroup}
\makeatother
\softneedspace{12\baselineskip}
\textbf{Re-weighting (TODO option 1.2): what it changes, and the pilot.}

\begin{explainbox}{How patient-equal re-weighting works, and what its percentages mean}
\emph{Unweighted (the stored bases).} Every cell has weight 1, so a patient's influence on the Stage~1 fit is its cell share (Figure~\ref{fig:02a}a, light bars): P09 28.5\%, P01 19.7\%, P02+P03+P04 8.2\%.

\emph{Patient-equal.} Cell $i$ of patient $p$ gets weight $w_i = 1/(13\, n_p)$, so each patient's weights sum to $1/13 = 7.7\%$ (the dashed line). P02+P03+P04 go from 8.2\% to $3/13 = 23.1\%$ of the fit; P09 goes from 28.5\% to 7.7\%. \emph{These percentages are shares of the total fit weight}, not of cells or of labels.

\emph{Mechanics.} Genes are re-standardized with the weighted mean and SD. PCA takes the eigenvectors of the weighted covariance $\Sigma_w = \sum_i w_i (x_i - \mu_w)(x_i - \mu_w)^\top$. FA has no sample weights in scikit-learn, so it is fitted to pseudo-data whose covariance equals $\Sigma_w$ (the FA likelihood depends on the data only through the covariance). A second arm subsamples 1,128 cells per patient (the smallest patient's count; 14,664 cells; 5 seeds). All 60,118 cells are then projected, and Stage~2 runs unchanged on the standard LOPO folds ($C$ tuned per fold).

\emph{What it does and does not change.} It raises the highest-prevalence patients' influence on Stage~1 and leaves Stage~2 untouched: 68.3\% of the malignant labels still come from them.
\end{explainbox}

\emph{Pilot, IFN-$\alpha$ only (still transductive).} Patient-equal weighting rotates PC8 (matched $|r|$ 0.90 at $K = 40$) more than FA7 (0.98), with the axis still in the subspace ($R^2$ 0.990, 0.994). Under L1 with $C$ tuned per fold there is no clear loss (pooled AP / median / P20 lift): PCA40 0.1213 / 9.25 / 2.69\pp{} vs 0.1153 / 8.91 / 1.42\pp{} for the unweighted transductive refit with the same code (the stored basis gives 0.1001 / 8.71 / 1.37\pp; section \emph{Option 2} explains the difference), median per-patient change $+1.35$\pp, no patient lower; FA40 0.0863 / 6.85 / 1.20 vs 0.0894 / 7.03 / 1.25\pp, median paired change $+0.07$\pp{} but all three summaries slightly lower. Under tuned (and fixed) L2, FA40 loses a median 1.13\pp{} per patient, with 5 of 7 patients lower. \src{2026-10-08\_d17-basis-variant-pilot; 2026-10-08\_team-report-devlog-pivot}

\emph{Targeting Stage~1 geometry or Stage~2 labels: the options} (W2.3). \textbf{You chose the Stage~2 labels on 2026-10-10}, for the runs with the new patient data; the Stage~1 options stay listed for reference.
\begin{itemize}
\item \emph{Stage~1 geometry} (who shapes the factors): weights in the DR fit (patient-equal, as piloted, or weights by cell type), equal-$n$ subsampling, healthy-cell downsampling before the fit, or an inductive or normal-only fit. Stage~2 stays as is. Read-out: the leverage metrics above, then held-out lift.
\item \emph{Stage~2 labels} (who the probe learns from): per-patient sample weights in the logistic loss (each patient's malignant cells sum to the same weight, likewise its normal cells), class balancing within each patient instead of pooled, capping or downsampling one patient's training cells, or choosing $C$ by patient-equal median lift instead of pooled AP. The basis stays as is. Read-out: per-patient, median and P20 lift.
\item Either changes the effective training data, so $C$ is retuned per condition, as in the pilot.
\item The evidence behind the choice: Stage~1 leverage of these patients is not beyond composition on IFN-$\alpha$, and on allograft and apoptosis it is beyond composition but not beyond that of other patients (D26, section~\ref{glance:02b}), while the Stage~2 labels are 68.3\% from three patients and lift ranks with prevalence (median within-representation Spearman $\rho = 0.821$ over 643 representations). \src{D6 in 2026-08-31\_d6-d9-meeting-followup; 2026-10-10\_d26-leverage-allograft-apoptosis-done}
\end{itemize}

\textbf{What this section means for healthy-cell downsampling} (your idea; the plan waits for the new cohort).
\begin{itemize}
\item \emph{In Stage~1, downsampling healthy cells is a re-weighting of the fit.} In the unweighted fit a patient's weight is its cell count. P09 holds 28.5\% of the cells and about the same share of the IFN-$\alpha$ $K = 40$ score energy (28.8\% PCA40, 27.5\% FA40): its cells are not unusually spread out, they are many. That is a statement about shares, not leverage: removing P09 moves the IFN-$\alpha$ axis more than all 20 of its composition-matched controls in 6 of 6 method $\times$ $K$ cells (Figure~\ref{fig:02b}c). Thinning P09's healthy cells in Stage~1 is the experiment that would show how much of the geometry its cell count sets. The patient with the most healthy cells per malignant cell, P01 (about 263), has the largest leverage of all (section~\ref{glance:02b}).
\item \emph{Energy per cell differs by cell type} (Figure~\ref{fig:02c}b: ncMono 3.25$\times$, HSC/MPP2 0.65$\times$). A cell-type-stratified downsampling keeps the Stage~1 energy budget only if it keeps the type proportions; changing them re-weights the fit toward or away from these types.
\item \emph{On IFN-$\alpha$ at $K = 40$, a Stage~1 gain should be small.} The malignant axis is carried by normal-cell variation: a basis fitted on normal cells only keeps PC8 (matched $|r|$ 0.9937) and FA7 (0.9937), so the axis is already in the subspace (on allograft PCA it is not: section~\ref{glance:02b}). Where a Stage~1 gain can show is at small $K$ and on panels where the malignant direction is not a leading factor: fewer healthy cells give the malignant cells more of the fit's variance and move that direction to an earlier factor.
\item \emph{In Stage~2 the ratio as such should already be neutralized} (an argument, not a result). With \texttt{class\_weight='balanced'} each class carries half of the loss however many healthy cells there are, so thinning healthy cells at random is the same objective at a smaller $C$, plus noise; and AP only ranks the held-out cells. With $C$ retuned, a Stage~2 gain would have to come from \emph{which} healthy cells are kept (for example thinning abundant types more than rare ones), not from the ratio itself. Your earlier donor downsampling (D10, fixed $C$) is consistent with this, since L2 held-out AP barely moved, but it did not retune $C$, so it is not a test (section \emph{Remaining TODOs}).
\item \emph{So, to explain a gain} (your main concern), the design separates the channels: downsample in Stage~1 only, in Stage~2 only, and in both, with $C$ retuned per condition and the held-out patient untouched (its prevalence $\pi$, the lift's baseline, does not move). Section \emph{Remaining TODOs}, box ``Why fewer healthy cells could help, and why it might not'', lists the four channels and the read-outs for each.
\end{itemize}
\src{2026-10-09\_report-reading-notes; 2026-10-09\_w8-healthy-downsampling-plan; 2026-10-08\_d17-basis-variant-pilot; 2026-10-08\_d17-composition-matched-controls; 2026-10-09\_d28-celltype-score-shares}

\begin{caveatbox}
\begin{itemize}
\item \textbf{One panel for leverage here} (IFN-$\alpha$; allograft rejection and apoptosis in section~\ref{glance:02b}; PCA and FA only, no FactoSig); score shares cover 41 panels. Score share is not leverage, and both use transductive bases.
\item \textbf{No significance}: 20 controls (rank-$p$ floor 0.048), 48 uncorrected contrasts per panel, 7 supported patients. On IFN-$\alpha$ the decision rule is a post-hoc reading (it was written for D26 after these results). A leave-out refit is an in-silico perturbation: it shows how much a patient's cells move an axis, not why. FA40 controls are bimodal, so FA40 ratios below 1 mean ``within control range'', not ``less leverage''.
\item \textbf{The P02+P03+P04 control is not free of them}: they supply 68.3\% of malignant cells, so the matched draws remove many of theirs too. The test isolates what is specific to their cells beyond the label count.
\item \textbf{The cell-type adjustment uses predicted labels} (label-transfer calls) and removes one confound; batch, sampling time and patient biology stay in the adjusted ratio. For P01 the reference matters (adjusted 1.03 or 1.08).
\item \textbf{P01's leverage is not explained}, and P02--P04 still dominate Stage~2 labels, selection and lift summaries.
\end{itemize}
\end{caveatbox}

\begin{teambox}
``We tested whether P02, P03 and P04 shape the factors in two ways. Their cells hold about their share of the factor variance on all 41 panels (1.04$\times$ after adjusting for cell type). On IFN-$\alpha$, refitting Stage~1 without them changes the malignant axis by under 1\% of its variance and not clearly more than removing the same number of random cells and malignant cells. On allograft rejection and apoptosis it moves 1.3 to 4 times as much as such random removals, mostly because of P03, but still by under 3\% of the axis's variance, and less than removing other patients, P01 above all. Where they clearly dominate is Stage~2: 68\% of the malignant labels.''

\textbf{Likely question:} ``So should we re-weight Stage~1 so every patient counts equally?''\\
\textbf{Answer:} Not as a fix for this: it would raise P02+P03+P04 from 8.2\% to 23.1\% of the fit (the pilot rotates PC8 to $|r|$ 0.90; no clear loss under tuned L1, but FA40 loses a median 1.13\pp{} under L2). If the worry is that three patients drive the results, the data place that mainly in Stage~2, where they supply 68\% of the labels, so re-weighting will target the Stage~2 labels in the runs with the new patients. The patient that moves the primary axis most is P01 (19.7\% of cells, 45 malignant), in 17 of the 18 method $\times$ $K$ cells of the three panels; its excess score energy is mostly matched by its cell-type mix, but its leverage is still unexplained.
\end{teambox}

\subsection*{Next step}
\begin{itemize}
\item \textbf{Done:} the allograft/apoptosis leverage test (D26, W2.4), which you ran in screen W2\_d26; verdict for PCA and FA: not supported. Section~\ref{glance:02b}, with its own next steps (W2.8 all panels, W2.9 P01, W2.2 FactoSig).
\item \textbf{Done:} the cell-type stratification of the score shares (D28, W1.9; Figure~\ref{fig:02c}).
\item Proposed: locating P02's and P07's excess that cell type does not explain, by factor and cell type (W1.10).
\item Re-weighting (W2.3): target decided on 2026-10-10 (the Stage~2 labels). The arm is designed with the runs on the new patient data, with $C$ retuned per condition; it runs inside the healthy-cell downsampling design as its Stage-2-only sample-weight scheme. Those runs leave out patients with no malignant cells and keep only patients with at least 100 malignant cells (W8.5; section \emph{Remaining TODOs}, subsection \emph{Proposed for the new data}). \src{2026-10-10\_decisions-new-run-inclusion-reweighting-lfs}
\end{itemize}

\subsection*{How this was computed}
{\small\raggedright
60,118 MRD CITE-seq cells (13 patients, 1,569 malignant); stored transductive Stage~1 bases, seed 42. \emph{Score shares (D16):} 41 biological panels $\times$ \{PCA, FA\} $\times$ $K \in \{5,10,20,40\}$, less 8 grid points where a panel has fewer genes than $K$: 320 bases with 5,870 factors (5,838 non-degenerate); energy on raw cohort-centred scores; no Stage~2 model, so no $C$ policy. \emph{Cell-type adjustment (D28):} the same 320 bases; 13 predicted cell types (\texttt{predicted.annotation}); expected energy = each cell's cell-type average over the cohort; adjusted ratio = observed / expected; nothing refit, no $C$ policy. \emph{Leverage (D17):} IFN-$\alpha$ (92 covered genes), PCA and FA at $K$ = 10, 20, 40, refit with the stored sklearn randomized solvers and genes standardized on the fit cells; reference = stored FA scores and, for PCA, the projection onto the stored loadings; size-matched controls seeds 1--3, composition-matched seeds 101--120; alignment only, no $C$ policy. \emph{D26} (allograft rejection, apoptosis): section~\ref{glance:02b}. \emph{Balanced pilot:} patient-equal weights over 13 patients; Stage~2 L1/L2 with $C$ tuned per outer fold by inner LOPO on the standard LOPO folds (training $n$ unchanged). Code: \texttt{scripts/\allowbreak mrd\_stage0\_2/\allowbreak workstreams/\allowbreak W1\_patient\_aggregation/\allowbreak run\_d16\_factor\_patient\_matrix.py}, \texttt{\ldots/\allowbreak W1\_patient\_aggregation/\allowbreak run\_d28\_celltype\_score\_shares.py}, \texttt{\ldots/\allowbreak W2\_stage1\_basis\_variants/\allowbreak run\_d17\_basis\_variant\_pilot.py}; figures and quoted ratios from this report's \texttt{figs\_02\_shape.py}. Artifacts (EXP-relative, under \texttt{analysis/\allowbreak scorecards/\allowbreak stage2\_diagnostic/}): \texttt{factor\_patient\_matrix\_v1/}, \texttt{factor\_patient\_matrix\_celltype\_v1/}, \texttt{basis\_variant\_pilot\_ifna\_v1/}. \src{2026-10-08\_d16-factor-patient-matrix; 2026-10-09\_d28-celltype-score-shares; 2026-10-08\_d17-basis-variant-pilot; 2026-10-08\_d17-composition-matched-controls; D6 in 2026-08-31\_d6-d9-meeting-followup; 2026-10-08\_team-report-devlog-pivot}\par}
